\nuevaReceta{Aceitunas en Ajedrea}{20 min (+ maceración)}{receta:aceitunas_ajedrea}

    \begin{minipage}[t]{0.35\textwidth}
        \section*{Ingredientes}
        \begin{itemize}[leftmargin=*]
            \item 1 kg de aceitunas verdes
            \item Agua
            \item Sal
            \item Ajedrea fresca o seca
            \item 4 Dientes de ajo
            \item 1 Limón
            \item 1 Hoja de laurel
            \item Hinojo (opcional)
        \end{itemize}
    \end{minipage}
    \hfill
    \begin{minipage}[t]{0.6\textwidth}
        \section*{Preparación}
        \begin{enumerate}
            \item \textbf{Preparar las aceitunas:} Lavar las aceitunas y partirlas con un golpe seco, procurando no romper el hueso.
            \item \textbf{Quitar el amargor:} Cubrirlas con agua limpia y cambiar el agua cada día durante 7-10 días, hasta que pierdan el amargor al gusto.
            \item \textbf{Preparar la salmuera:} Disolver aproximadamente 70 g de sal por cada litro de agua. Remover hasta que la sal quede completamente integrada.
            \item \textbf{Preparar el aliño:} Machacar ligeramente los ajos y cortar el limón en trozos o rodajas gruesas.
            \item \textbf{Llenar el recipiente:} Colocar las aceitunas en un tarro limpio, alternándolas con la ajedrea, los ajos, el limón, el laurel y, si se desea, un poco de hinojo.
            \item \textbf{Cubrir:} Verter la salmuera hasta que todas las aceitunas queden completamente sumergidas.
            \item \textbf{Macerar:} Cerrar el recipiente y dejar reposar en un lugar fresco y oscuro durante al menos 2 semanas antes de probarlas.
        \end{enumerate}
    \end{minipage}

    \vspace{1em}
    \begin{tcolorbox}[colback=orange!5, colframe=orange!50, title=Consejo, size=small]
        Mantén siempre las aceitunas cubiertas por la salmuera y utiliza utensilios limpios para sacarlas. El sabor de la ajedrea se intensifica con los días de reposo.
    \end{tcolorbox}
\end{receta}
